% ==============================================================================
% PLANTILLA MAESTRA - UnADM (instrucciones comunes)
% ==============================================================================
% Uso rapido:
% 1) Duplica la plantilla y renombra la copia por semana/actividad.
% 2) Actualiza \documenttitle, \documentsubtitle y \documentsubject.
% 3) Lee la planeacion y NO pegues las instrucciones completas dentro del PDF.
%    Usa la planeacion como lista de cumplimiento y redacta con criterio propio.
% 4) Elabora productos visuales preferentemente en LaTeX/TikZ para mantener
%    nitidez y evidencia editable; si usas herramientas externas, inserta
%    capturas de alta resolucion y conserva evidencia dentro del PDF.
% 5) Entrega en PDF con la nomenclatura solicitada en el aula.
%
% Formato institucional (recomendado):
% - Fuente: Helvetica / sans-serif equivalente (\usepackage{helvet}).
% - Interlineado: 1.5 (\onehalfspacing).
% - Citas: autor-año mediante natbib (\citep, \citet, \citep[e.g.][]{clave}).
% - Bibliografía: archivo .bib con formato APA 7; todas las citas del texto
%   deben aparecer en la bibliografía y viceversa.
% - Estructura mínima: portada, resumen/abstract, desarrollo, conclusión y
%   bibliografía.
%
% Reglas para productos visuales and tablas extensas:
% - Para tablas o cuadros amplios use longtable + booktabs y evite líneas
%   verticales; controlar espaciado con \setlength{\tabcolsep} y
%   \renewcommand{\arraystretch}.
% - Para tablas en página propia: use \clearpage, \pagestyle{empty},
%   entorno landscape, reducir tamaño con \scriptsize y luego restaurar
%   \pagestyle{fancy} al finalizar la tabla.
%
% Buenas prácticas de redacción académica (guía operativa):
% - No escribir para llenar espacio; escribir para sostener y defender una idea.
% - Cada concepto debe responder: qué es, qué características tiene, cómo
%   funciona, qué problema resuelve, qué límite presenta y qué implicación
%   jurídica posee.
% - Evitar listas sin análisis: tras cada dato relevante incluya su consecuencia
%   o aplicación jurídica.
% - Primera persona: usar moderadamente ("Considero que", "Desde esta
%   perspectiva"); evite expresiones coloquiales como "yo creo".
%
% Protocolos de integridad y uso de IA:
% - Si se usa IA, el resultado debe ser leído, corregido y transformado con
%   criterio propio. No entregar texto sin procesamiento.
% - Incluir siempre juicio propio y referencias que sostengan la postura.
%
% Checklist operativo antes de entregar:
% [ ] El producto coincide con la técnica solicitada en la planeación.
% [ ] El objetivo y el contenido están alineados con la planeación.
% [ ] El producto visual está rotulado y presentado como se pide.
% [ ] Hay citas APA autor-año dentro del texto y bibliografía completa.
% [ ] No se incluyeron instrucciones copiadas de la planeación dentro del
%     contenido entregado.
% [ ] Ortografía y redacción revisadas; la argumentación es propia.
% ==============================================================================

% CREACION DEL DOCUMENTO
\documentclass[
	spanish,
	letterpaper, oneside
]{article}

% ==============================================================================
% INFORMACION DEL DOCUMENTO
% Actualizar en cada actividad.
% Ejemplos:
% \def\documenttitle {Cuadro comparativo de criterios de redaccion academica}
% \def\documentsubtitle {Actividad X - Redaccion en contextos virtuales}
% \def\documentsubject {Actividad Semana X}
% ==============================================================================
\def\documenttitle {Portafolio de evidencias de redaccion en contextos virtuales}
\def\documentsubtitle {Integracion de actividades 1 a 6}
\def\documentsubject {Actividad Semana 9}

\def\documentauthor {Martin Jonathan de la Cruz}
\def\coursename {Redaccion en contextos virtuales}
\def\coursecode {017}

\def\universityname {Universidad Abierta y a Distancia de Mexico}
\def\universityfaculty {Licenciatura en Derecho}
\def\universitydepartment {Redaccion en contextos virtuales}
\def\universitydepartmentimage {departamentos/UnADM}
\def\universitydepartmentimagecfg {height=1.57cm}
\def\universitylocation {Roma Norte, Ciudad de Mexico}

% INTEGRANTES, PROFESORES Y FECHAS
\def\authortable {
	\begin{tabular}{ll}
		\textbf{Alumno:}
		& \begin{tabular}[t]{l}
			 Martin Jonathan de la Cruz \\
		\end{tabular} \\
		Matricula: & ES2611202040 \\

		\textit{Profesora:}
		& \begin{tabular}[t]{l}
			Ana Rosa \\ Arceo Luna
		\end{tabular} \\
		& \\
		\multicolumn{2}{l}{Fecha de realizacion: \today} \\
		\multicolumn{2}{l}{\universitylocation}
	\end{tabular}
}

% IMPORTACION DEL TEMPLATE
\input{template}

% FORMATO SOLICITADO / USADO EN ACTIVIDADES CORREGIDAS
\usepackage{helvet}
\renewcommand{\familydefault}{\sfdefault}

% TABLAS Y PRODUCTOS VISUALES
\usepackage{array}
\usepackage{longtable}
\usepackage{booktabs}
\usepackage{pdflscape}
\usepackage{tikz}
\usetikzlibrary{
	arrows.meta,
	positioning,
	calc,
	matrix,
	fit,
	shapes.geometric,
	shadows.blur
}

% CITAS EN FORMATO AUTOR-ANIO, COMPATIBLE CON APA 7 EN EL CUERPO DEL TEXTO
% Usar:
% - Cita parentetica: \citep{clave}
% - Cita narrativa: \citet{clave}
% - Varias fuentes: \citep{clave1,clave2}
\setcitestyle{authoryear,open={(},close={)}}

% ==============================================================================
% MARCA DE AGUA EN PORTADA
% - Se aplica solo a la portada porque usa \AddToShipoutPictureBG*.
% - Para apagarla en alguna actividad: \def\coverwatermarkenabled {false}
% - Usar opacidad baja para que no compita con titulo, datos ni logotipo.
% ==============================================================================
\def\coverwatermarkenabled {true}
\def\coverwatermarkimage {img/departamentos/UnADM.pdf}
\def\coverwatermarkopacity {0.16}
\def\coverwatermarkwidth {0.72\paperwidth}
\def\coverwatermarkxshift {0cm}
\def\coverwatermarkyshift {-0.35cm}

\newcommand{\insertcoverwatermark}{%
	\ifthenelse{\equal{\coverwatermarkenabled}{true}}{%
		\AddToShipoutPictureBG*{%
			\begin{tikzpicture}[remember picture,overlay]
				\node[
					opacity=\coverwatermarkopacity,
					inner sep=0pt,
					xshift=\coverwatermarkxshift,
					yshift=\coverwatermarkyshift
				] at (current page.center) {%
					\includegraphics[width=\coverwatermarkwidth]{\coverwatermarkimage}%
				};
			\end{tikzpicture}%
		}%
	}{}%
}

% ==============================================================================
% AYUDAS DE REDACCION
% Quitar o comentar \pendiente cuando el documento final este terminado.
% ==============================================================================
\newcommand{\pendiente}[1]{\textcolor{red}{[PENDIENTE: #1]}}

% RESUMENES INTERNOS DEL PORTAFOLIO
% A diferencia de abstractd, este entorno no reinicia secciones ni subsecciones.
\newenvironment{portafolioabstract}{%
	\subsubsection*{\nameabstract}%
}{%
	\par
}

% CAJAS USADAS EN LA EVIDENCIA ORIGINAL DE LA ACTIVIDAD 5
\newcommand{\portafolioexercisebox}[2]{%
	\par\medskip
	\noindent\fcolorbox{gray!55}{gray!25}{%
		\parbox{\dimexpr\linewidth-2\fboxsep-2\fboxrule\relax}{\textbf{#1}}%
	}%
	\par\noindent
	\fcolorbox{gray!55}{gray!10}{%
		\parbox{\dimexpr\linewidth-2\fboxsep-2\fboxrule\relax}{#2}%
	}%
	\par\medskip
}

\begin{document}

% PORTADA
\insertcoverwatermark
\templatePortrait

% CONFIGURACION DE PAGINA Y ENCABEZADOS
\templatePagecfg
\onehalfspacing

% ==============================================================================
% RESUMEN / ABSTRACT
% Debe ser breve. Formula recomendada:
% Tema + objetivo + tecnica/producto + enfoque personal.
% ==============================================================================
\begin{abstractd}
	\textbf{Objetivo:} Integrar y analizar las actividades 1 a 6 de la asignatura para mostrar, mediante un portafolio de evidencias, el proceso de desarrollo de competencias de comunicacion y redaccion academica en entornos virtuales.

	\textbf{Producto elaborado:} Portafolio de evidencias con estructura propia de actividad academica que recupera el estudio de caso, el cuadro comparativo, el juego de roles, la lista de verificacion, el ejercicio de transformacion de texto y el mapa conceptual elaborados a lo largo del curso, e incorpora cada uno de esos trabajos en su version completa dentro del desarrollo.

	\textbf{Enfoque de analisis:} El conjunto de evidencias permite observar una progresion: primero se identifica el problema comunicativo, despues se aplican criterios de claridad y formalidad, luego se revisa la legibilidad del texto propio y, finalmente, se sistematiza la relacion entre autonomia, organizacion y calidad de la escritura academica \citep{ayalaTallerLecturaRedaccion2010,nunezcortesEscrituraAcademicaTeoria2015,vitoriaEscrituraAcademicaFormacion2019}.
\end{abstractd}

% TABLA DE CONTENIDOS - INDICE
\templateIndex

% CONFIGURACIONES FINALES
\templateFinalcfg

% ======================= INICIO DEL DOCUMENTO =======================

\section{Introduccion}

La escritura academica en entornos virtuales no se desarrolla en un solo ejercicio; se construye por acumulacion, revision y transferencia de criterios entre actividades distintas. Por ello, el portafolio de evidencias no debe entenderse como un archivo pasivo de trabajos entregados, sino como un instrumento de lectura retrospectiva que permite reconocer avances, limites y relaciones entre productos que, en apariencia, fueron resueltos por separado.

En esta asignatura, las actividades 1 a 6 abordaron problemas complementarios: el proceso de comunicacion, los criterios de la redaccion academica, la escritura laboral, la evaluacion de textos digitales, la transformacion de registros coloquiales y la organizacion conceptual del aprendizaje autonomo. Aunque la planeacion contempla integrar tres evidencias, en este documento se recuperan las seis actividades para mostrar de manera completa la trayectoria formativa y el modo en que cada producto aporta a una misma competencia: comunicar con claridad, coherencia, cohesion y pertinencia en contextos virtuales \citep{correaManualBasicoEscritura2018,nunezcortesEscrituraAcademicaTeoria2015}.

Desde esta perspectiva, el portafolio funciona como una sintesis argumentada del proceso. No solo reune productos terminados; tambien explica que problema resolvio cada uno, que criterio fortalecio y de que manera contribuyo a una comprension mas rigurosa de la redaccion academica.

\section{Desarrollo}

El desarrollo no presenta seis trabajos como piezas aisladas, sino como la evidencia de una competencia que se construye paso a paso. El recorrido comienza con un conflicto comunicativo, avanza hacia la identificacion de criterios de calidad textual y culmina con la capacidad de transformar, evaluar y organizar la escritura academica dentro de un entorno virtual.

Cada actividad se incorpora en su version completa para hacer visible que aprendizaje produjo y como se relaciona con las demas. La Actividad 1 analiza el efecto de un mensaje mal formulado; la Actividad 2 establece criterios para distinguir una redaccion espontanea de una academica; la Actividad 3 traslada esos criterios al ambito laboral; la Actividad 4 examina la legibilidad de un texto digital; la Actividad 5 demuestra el proceso de reescritura; y la Actividad 6 organiza conceptualmente la relacion entre autonomia, responsabilidad y calidad textual.

En conjunto, estas evidencias muestran una transformacion concreta: pasar de escribir como reaccion inmediata a escribir con intencion, estructura y conciencia del receptor. Esa trayectoria constituye el eje del portafolio y permite observar el aprendizaje directamente en los productos elaborados.

\subsection{Actividad 1. Estudio de caso: elementos del proceso de comunicacion}
\input{anexos-portafolio/actividad-1}

\clearpage
\subsection{Actividad 2. Cuadro comparativo de criterios de la redaccion academica}
\input{anexos-portafolio/actividad-2}

\clearpage
\subsection{Actividad 3. Comunicacion efectiva en el ambito laboral: juego de roles}
\begin{portafolioabstract}
	\textbf{Objetivo:} Aplicar los principios de la comunicación efectiva en el ámbito laboral mediante el análisis de correos electrónicos y la redacción de mensajes adecuados al receptor, la intención comunicativa y el nivel de formalidad requerido.

	\textbf{Producto elaborado:} Juego de roles con análisis de ejemplos, redacción de un correo inadecuado, reformulación formal y reflexión final.

	\textbf{Enfoque de análisis:} La comunicación laboral escrita no se limita a transmitir información; también ordena responsabilidades, cuida la relación profesional y reduce el margen de malentendido.
\end{portafolioabstract}

\subsubsection{Introducción}

La comunicación escrita en el ámbito laboral exige más que corrección gramatical. Un correo puede contener una solicitud válida y, aun así, producir resistencia si el tono es impreciso, apresurado o descortés. En contextos profesionales, la forma del mensaje también comunica: muestra respeto por el receptor, claridad sobre la tarea y conciencia de las consecuencias que una petición puede generar.

Desde esta perspectiva, la redacción académica y profesional funciona como una herramienta de control del sentido. No se trata de escribir de manera rígida, sino de elegir palabras, estructura y nivel de formalidad de acuerdo con el contexto. La lectura y la redacción deben entenderse como procesos vinculados: primero se interpreta la situación comunicativa y después se construye un mensaje adecuado a su propósito \citep{ayalaTallerLecturaRedaccion2010}. Como señalan los estudios sobre escritura académica, producir un texto implica organizar ideas, adecuarlas a una situación comunicativa y hacerlas comprensibles para un lector concreto \citep{nunezcortesEscrituraAcademicaTeoria2015,vitoriaEscrituraAcademicaFormacion2019}. Por ello, esta actividad analiza dos correos laborales, identifica errores y aciertos, y propone una versión inadecuada y una versión formal a partir de una situación de trabajo.

\subsubsection{Análisis de ejemplos}

Los dos correos de la planeación muestran una diferencia central: el primero comunica urgencia, pero no construye condiciones para la colaboración; el segundo conserva la urgencia, pero la organiza con cortesía, estructura y precisión. En términos de redacción, el problema no es pedir algo con rapidez, sino hacerlo sin asunto claro, sin contexto suficiente y con expresiones que pueden ser interpretadas como presión o descuido profesional.

\begin{longtable}{@{}p{0.22\linewidth}p{0.34\linewidth}p{0.34\linewidth}@{}}
\caption{Análisis comparativo de los correos de ejemplo}\label{tab:analisis-correos}\\
\toprule
\textbf{Criterio} & \textbf{Correo inadecuado} & \textbf{Correo adecuado} \\
\midrule
\endfirsthead
\toprule
\textbf{Criterio} & \textbf{Correo inadecuado} & \textbf{Correo adecuado} \\
\midrule
\endhead
Asunto & El asunto ``duda'' es demasiado general y no permite anticipar la finalidad del mensaje. En un entorno laboral, un asunto impreciso obliga al receptor a interpretar antes de actuar. & El asunto ``Solicitud de envío de reporte pendiente'' indica con claridad la acción esperada y el documento requerido. Facilita priorizar, localizar y responder el mensaje. \\
\addlinespace
Tono & Usa expresiones coloquiales como ``oye'', ``qué onda'' y ``urge''. El mensaje puede percibirse como exigente, aunque la intención sea resolver una necesidad inmediata. & Mantiene cortesía mediante saludo, agradecimiento y cierre formal. La solicitud es directa, pero no invasiva. \\
\addlinespace
Claridad & La petición no explica con precisión qué documento se necesita, para qué se requiere ni por qué el tiempo es relevante. & Explica que se da seguimiento a una solicitud previa, identifica el reporte mensual y justifica la necesidad por una entrega programada. \\
\addlinespace
Estructura & Carece de apertura formal, contexto, desarrollo y cierre profesional. La información aparece como presión acumulada. & Presenta saludo, contexto, solicitud, disponibilidad para aclaraciones, agradecimiento y firma. Esa estructura reduce ambigüedad. \\
\addlinespace
Adecuación al receptor & No considera la posición ni el tiempo del destinatario. La urgencia se traslada al receptor sin construir una solicitud profesional. & Reconoce al destinatario con tratamiento formal y deja abierta la posibilidad de comentarios o información adicional. \\
\bottomrule
\end{longtable}

En el primer correo se identifican al menos tres errores relevantes: asunto vago, tono informal y falta de estructura. También se observa una urgencia mal gestionada, porque el mensaje no explica el motivo de la premura ni ofrece información suficiente para que la otra persona responda con eficiencia. En cambio, el segundo correo presenta tres aciertos claros: asunto específico, cortesía profesional y solicitud contextualizada. Esta diferencia confirma que la claridad no depende solo de usar menos palabras, sino de ordenar la información para que el receptor entienda qué se solicita, por qué se solicita y cómo puede responder.

\subsubsection{Juego de roles}

Para el juego de roles se eligió la situación laboral: \textbf{dar seguimiento a una actividad pendiente}. El supuesto es el siguiente: una persona integrante de un equipo jurídico-administrativo necesita solicitar a una compañera el envío de observaciones pendientes sobre un documento que debe integrarse a un expediente de trabajo antes del cierre del día.

\paragraph{Versión 1: correo inadecuado}

\begin{quote}
\textbf{Asunto:} pendientes

Hola, Laura:

Oye, sigo esperando tus observaciones del documento. Me dijiste que me las ibas a mandar y todavía no me llega nada. La verdad sí me urge porque tengo que cerrar esto hoy y no puedo avanzar si no me contestas.

Mándamelo en cuanto puedas porque ya vamos tarde.

Gracias.
\end{quote}

Esta versión es coherente con una situación laboral real, pero contiene errores típicos: el asunto es impreciso, el tono transmite molestia, no se identifica con claridad el documento, no se explica el propósito institucional de la solicitud y la urgencia se expresa como presión personal. El mensaje podría generar una respuesta defensiva porque coloca el problema sobre la destinataria sin organizar la petición de manera profesional.

\paragraph{Versión 2: correo formal y adecuado}

\begin{quote}
\textbf{Asunto:} Seguimiento a observaciones pendientes del documento de expediente

Estimada Laura:

Recibe un cordial saludo.

Me permito dar seguimiento a las observaciones pendientes sobre el documento de integración del expediente que revisamos durante la mañana. Agradecería mucho si pudieras compartirme tus comentarios antes de las 16:00 horas, ya que son necesarios para consolidar la versión final y enviarla dentro del plazo programado.

Quedo atento a cualquier duda o ajuste que consideres necesario revisar antes del envío.

Agradezco de antemano tu apoyo.

Atentamente,

Martín Jonathan de la Cruz
\end{quote}

La segunda versión conserva la misma intención comunicativa, pero modifica la relación con el receptor. En lugar de trasladar presión, organiza la solicitud: identifica el documento, precisa el plazo, justifica la urgencia y mantiene apertura para resolver dudas. Desde una lógica profesional, este tipo de redacción no solo mejora la cortesía; también aumenta la posibilidad de obtener una respuesta útil, porque el receptor recibe instrucciones claras y un contexto suficiente para actuar.

\subsubsection{Reflexión final}

Entre ambas versiones realicé cambios de tono, estructura y precisión. El correo inadecuado parte de la urgencia personal y se expresa con informalidad; el correo formal parte del contexto laboral y convierte la urgencia en una solicitud organizada. Apliqué tres principios de comunicación efectiva: claridad, porque el mensaje indica qué documento se requiere y en qué plazo; formalidad, porque el saludo, el tratamiento y el cierre corresponden al ámbito profesional; y cortesía, porque la petición reconoce al receptor y deja abierta la posibilidad de aclaraciones. Considero que esta diferencia es importante para la formación jurídica, ya que un mensaje mal redactado puede afectar la coordinación, generar tensión innecesaria o debilitar la confianza en quien lo emite. La escritura profesional debe servir para resolver, no para aumentar el ruido del problema.

\subsubsection{Conclusión}

El ejercicio permite observar que la comunicación laboral escrita requiere equilibrio entre intención, receptor y contexto. Un mensaje puede ser urgente sin ser brusco, directo sin ser descortés y breve sin ser ambiguo. La diferencia está en la forma de organizar la información: un asunto claro, una apertura respetuosa, una solicitud precisa, una justificación suficiente y un cierre profesional convierten una petición cotidiana en un acto comunicativo eficaz.

Desde mi perspectiva, la principal enseñanza de la actividad es que escribir bien en el ámbito profesional significa anticipar la interpretación del receptor. La redacción no solo expresa lo que se necesita; también construye la relación desde la cual se espera una respuesta. En ese sentido, la claridad, la formalidad y la cortesía no son adornos del correo laboral, sino condiciones mínimas para trabajar con responsabilidad en entornos académicos, administrativos y jurídicos.


\clearpage
\subsection{Actividad 4. Lista de verificacion para evaluar textos digitales}
\input{anexos-portafolio/actividad-4}

\clearpage
\subsection{Actividad 5. Ejercicio de transformacion: de redaccion libre a redaccion academica}
\begin{portafolioabstract}
	\textbf{Objetivo:} Analizar la diferencia entre redacción libre y redacción académica mediante la reescritura de un texto coloquial sobre el estudio en línea.

	\textbf{Producto elaborado:} Presentación del texto base, identificación de sus rasgos de oralidad e imprecisión y elaboración de una versión académica reformulada en 150--200 palabras.

	\textbf{Enfoque de análisis:} La escritura universitaria no se limita a corregir formas superficiales; exige organizar ideas, jerarquizar argumentos y convertir una experiencia personal en un mensaje claro, interpretable y pertinente para un contexto formativo \citep{nunezcortesEscrituraAcademicaTeoria2015,correaManualBasicoEscritura2018}.
\end{portafolioabstract}

\subsubsection{Introducción}

En redacción académica, la forma de decir una idea es tan importante como la idea misma. Un texto puede comunicar una experiencia válida, pero si está redactado con coloquialismos, saltos lógicos o imprecisión léxica, pierde fuerza argumentativa en contextos universitarios.

Ese es el problema que atiende esta actividad: transformar un mensaje cotidiano sobre el estudio en línea en un texto académico que conserve el significado, pero mejore su estructura, su formalidad y su coherencia. La redacción académica exige claridad, intención comunicativa y articulación de ideas mediante conectores, no solo sustitución de palabras informales \citep{correaManualBasicoEscritura2018,vitoriaEscrituraAcademicaFormacion2019}.

Desde esta perspectiva, el ejercicio de transformación permite observar cómo cambia la eficacia comunicativa cuando el texto pasa de una redacción libre a una redacción orientada por criterios académicos.

\subsubsection{Transformación de redacción libre a redacción académica}

La transformación de redacción libre a redacción académica exige pasar de una expresión espontánea a una construcción discursiva organizada, precisa y orientada por una intención comunicativa explícita. En este marco, escribir no significa solo trasladar ideas, sino jerarquizarlas, vincularlas mediante relaciones lógicas y presentarlas con un registro pertinente para el contexto universitario. Desde esta perspectiva, la diferencia entre ambas formas de escritura no radica en la eliminación de la voz personal, sino en el nivel de control textual que permite convertir una experiencia cotidiana en un argumento claro y comprensible \citep{nunezcortesEscrituraAcademicaTeoria2015,correaManualBasicoEscritura2018,vitoriaEscrituraAcademicaFormacion2019}.

Con base en ello, el ejercicio consiste en transformar un párrafo originalmente redactado en estilo coloquial hacia una versión académica que conserve la idea central, pero mejore su coherencia, su precisión y su fuerza explicativa. El párrafo original es el siguiente:

\portafolioexercisebox{Texto original}{%
``Hoy en día estudiar en línea está medio complicado porque la neta te distraes mucho. O sea, estás leyendo y de repente ya estás viendo redes o contestando mensajes. Además, siento que no entiendes igual que cuando lees en papel, porque como que solo lees rápido y ya. Pero bueno, también es más práctico porque puedes leer en cualquier lado.''%
}

El proceso de transformación en lenguaje académico se desarrolló al identificar la idea central del texto, eliminar marcas de oralidad, reorganizar la secuencia argumentativa y establecer conectores que articularan de manera coherente causas, efectos y valoraciones. También se ajustó el vocabulario para reducir ambigüedades y se revisó la progresión temática con el fin de mantener unidad y claridad en todo el párrafo.

\portafolioexercisebox{Versión académica}{%
El estudio en línea presenta retos vinculados con la gestión de la atención y la calidad de la comprensión lectora. En los entornos digitales, la lectura compite de manera constante con notificaciones, redes sociales y otros estímulos que fragmentan la concentración. Esta interferencia no solo interrumpe el ritmo de trabajo, sino que también puede reducir la profundidad con la que se procesa la información.\par

Además, la lectura en pantalla suele realizarse con mayor rapidez y con menos pausas de reflexión que la lectura en formato impreso, especialmente cuando la persona estudiante no aplica estrategias de autorregulación. Como consecuencia, parte del contenido puede revisarse de manera superficial, lo que afecta la retención y limita el análisis posterior.\par

No obstante, el formato digital ofrece una ventaja operativa clara: permite acceder a materiales desde distintos espacios y dispositivos. Por ello, su eficacia académica no depende únicamente del soporte, sino de las prácticas de lectura y concentración que desarrolle el estudiante. En este sentido, fortalecer hábitos de atención, lectura crítica y manejo del tiempo resulta esencial para aprovechar de manera efectiva la flexibilidad propia de la educación en línea.%
}

Como puede apreciarse en el resultado, la versión académica conserva el contenido esencial del párrafo original, pero mejora su estructura lógica, su formalidad y su capacidad explicativa. La información deja de presentarse como una impresión dispersa y se convierte en una argumentación ordenada, con relaciones explícitas entre problema, consecuencias y valoración final.

\subsubsection{Conclusión}

El ejercicio de transformación permitió comprobar que la escritura académica no se limita a corregir palabras informales, sino a reorganizar el pensamiento para hacerlo más claro, preciso y argumentativamente consistente. La versión original comunicaba una experiencia auténtica, pero su estructura reducía la fuerza del mensaje y dificultaba su proyección en un contexto universitario.

La versión reformulada mantuvo la idea central y mejoró su legibilidad, su coherencia y su capacidad explicativa. En conclusión, transformar redacción libre en redacción académica significa pasar de una expresión espontánea a una construcción discursiva con intención, control y sentido formativo. Ese es, en última instancia, el valor real del ejercicio.


\clearpage
\subsection{Actividad 6. Mapa conceptual sobre autogestion del aprendizaje y redaccion academica}
\begin{portafolioabstract}
	\textbf{Objetivo:} Analizar cómo la educación superior abierta y a distancia desplaza parte del control del aprendizaje hacia el estudiante y convierte la redacción académica en una evidencia central de comprensión, seguimiento y responsabilidad.

	\textbf{Producto elaborado:} Mapa conceptual que organiza, jerarquiza y relaciona conceptos clave del bloque temático mediante nodos, conectores y palabras enlace con función explicativa.

	\textbf{Enfoque de análisis:} Se sostiene que la calidad de la escritura en entornos virtuales no depende solo del manejo formal de la lengua, sino de un sistema compuesto por autogestión, organización del tiempo, responsabilidad académica y criterios de claridad, coherencia y cohesión \citep{nunezcortesEscrituraAcademicaTeoria2015,correaManualBasicoEscritura2018,vitoriaEscrituraAcademicaFormacion2019}.
\end{portafolioabstract}

\subsubsection{Introducción}

En la educación superior abierta y a distancia, escribir bien no es un adorno académico, sino una condición de funcionamiento del aprendizaje. Cuando la mayor parte de la interacción ocurre mediante plataformas, mensajes, consignas y entregables escritos, la redacción deja de ser una habilidad secundaria y se convierte en el medio principal para demostrar comprensión, seguimiento y criterio.

La actividad parte de un problema concreto: explicar por qué la calidad de la escritura en entornos virtuales no depende solo del dominio formal de la lengua, sino de la relación entre autonomía, organización del tiempo, responsabilidad académica y control del propio proceso de estudio. Si esa estructura no se comprende, la escritura se reduce a respuestas fragmentarias, la comunicación pierde eficacia y el aprendizaje deja de ser visible como proceso riguroso \citep{correaManualBasicoEscritura2018,nunezcortesEscrituraAcademicaTeoria2015}.

Desde esta perspectiva, el mapa conceptual funciona como un instrumento de diagnóstico y no solo como un recurso visual. Su utilidad consiste en mostrar que la redacción académica opera como resultado de un sistema: la modalidad virtual exige autogestión; la autogestión necesita disciplina; y ambas condiciones terminan reflejándose en la claridad, coherencia y pertinencia de lo que se comunica.

\subsubsection{Sistema conceptual de la redacción académica en entornos virtuales}

\paragraph{Contexto de análisis}

La educación superior abierta y a distancia redistribuye una parte sustantiva del control del aprendizaje hacia el estudiante. Esa reconfiguración no solo cambia la forma de estudiar; también modifica la forma en que se administra el tiempo, se procesan las instrucciones y se comunica el conocimiento. En este modelo, la escritura funciona como registro operativo del aprendizaje: deja evidencia de comprensión, seguimiento, disciplina y capacidad de respuesta.

Por eso, el problema no se agota en redactar sin errores formales. El punto de fondo consiste en sostener un proceso de estudio capaz de producir textos claros, ordenados y pertinentes dentro de una modalidad flexible, mediada por plataformas y marcada por una menor supervisión inmediata. En otras palabras, la calidad textual no puede separarse de la capacidad de autogestión, porque ambas forman parte del mismo sistema de desempeño académico \citep{vitoriaEscrituraAcademicaFormacion2019,correaManualBasicoEscritura2018}.

\paragraph{Conceptos fundamentales}

\textbf{Educación abierta y a distancia}. Es una modalidad caracterizada por flexibilidad espacial y temporal. Su valor no se encuentra solo en ampliar el acceso, sino en exigir nuevas formas de organización, seguimiento y permanencia dentro del proceso formativo \citep{vitoriaEscrituraAcademicaFormacion2019}.

\textbf{Autogestión del aprendizaje}. Consiste en planificar, supervisar y ajustar el propio estudio. En términos prácticos, funciona como un sistema de control personal: permite administrar tiempos, recursos y estrategias para sostener continuidad y evitar dispersión.

\textbf{Organización del tiempo}. Es la condición operativa que convierte la flexibilidad en avance real. Cuando el tiempo no se estructura, la modalidad virtual pierde eficiencia; cuando se planifica, el aprendizaje adquiere ritmo, secuencia y trazabilidad.

\textbf{Responsabilidad académica}. Implica cumplir compromisos, revisar indicaciones, sostener lectura y cuidar la calidad de lo que se entrega. No es un rasgo abstracto, sino una práctica verificable en la constancia del trabajo y en el nivel de precisión del texto.

\textbf{Aprendizaje autónomo}. Se refiere a la capacidad de avanzar con iniciativa propia, resolver dudas, corregir errores y mantener continuidad sin depender de una vigilancia constante. Esta autonomía no significa aislamiento, sino participación activa y consciente en la propia formación \citep{correaManualBasicoEscritura2018}.

\textbf{Redacción académica}. Es una práctica de comunicación que organiza ideas con orden lógico, lenguaje pertinente y propósito claro. Su función no es adornar el contenido, sino volverlo interpretable, defendible y útil para una comunidad académica \citep{nunezcortesEscrituraAcademicaTeoria2015}.

\textbf{Claridad}. Permite que la idea principal sea comprensible sin ambigüedades innecesarias. Un texto claro reduce ruido, orienta la lectura y facilita la interpretación correcta del mensaje.

\textbf{Coherencia}. Asegura unidad temática y continuidad del argumento. Gracias a ella, cada parte del texto responde a una misma línea de sentido y evita contradicciones internas.

\textbf{Cohesión}. Conecta las partes del discurso mediante relaciones explícitas. Su función es enlazar oraciones, párrafos y conceptos para que el texto avance como sistema y no como fragmentos aislados.

\textbf{Formalidad}. Ajusta el registro al contexto universitario. Esto implica precisión léxica, moderación expresiva y respeto por las convenciones propias de la escritura académica.

\textbf{Comunicación en entornos virtuales}. Es el espacio donde convergen autonomía, responsabilidad y redacción. En esta modalidad, gran parte de la presencia académica se construye a través de mensajes, participaciones y productos escritos; por eso, la calidad de la comunicación depende directamente de la calidad del proceso de estudio.

\paragraph{Síntesis del mapa conceptual}

La actividad original representó estas relaciones mediante un mapa conceptual elaborado en TikZ. Su aporte principal consistió en mostrar que la modalidad virtual impone condiciones, la autogestión traduce esas condiciones en acciones concretas y la redacción académica aparece como la salida verificable de ese proceso. No se trató de una lista de conceptos aislados, sino de una estructura donde cada conector verbal expresaba dependencia, consecuencia o criterio de valoración.

En términos de contenido, el mapa permitió observar tres relaciones decisivas. Primero, que la educación abierta y a distancia exige flexibilidad, pero también demanda organización del tiempo y aprendizaje autónomo. Segundo, que la responsabilidad académica se manifiesta en prácticas concretas como lectura, revisión, corrección y cumplimiento. Tercero, que la calidad textual se valora por criterios como claridad, coherencia, cohesión y formalidad, todos ellos vinculados con la forma en que el estudiante gestiona su propio proceso de estudio.

\clearpage
\begingroup
\pagestyle{empty}
\thispagestyle{empty}
\begin{landscape}
\begin{figure}[p]
\centering
\vspace*{-0.25cm}
\resizebox{0.96\linewidth}{!}{%
\begin{tikzpicture}[
	x=1cm,
	y=1cm,
	concept/.style={
		rectangle,
		rounded corners=3pt,
		draw=dkgray,
		line width=0.45pt,
		fill=lgray,
		align=center,
		text width=2.85cm,
		minimum height=0.85cm,
		inner sep=4pt,
		font=\sffamily\footnotesize
	},
	root/.style={
		concept,
		fill=cardinalred!14,
		draw=cardinalred!70!black,
		text width=4.25cm,
		minimum height=1.25cm,
		font=\bfseries\sffamily\normalsize
	},
	branch/.style={
		concept,
		fill=dkcyan!10,
		draw=dkcyan!60!black,
		text width=2.75cm,
		minimum height=0.95cm,
		font=\bfseries\sffamily\footnotesize
	},
	subbranch/.style={
		rectangle,
		rounded corners=2pt,
		draw=dkgray!75,
		line width=0.35pt,
		fill=white,
		align=center,
		text width=2.55cm,
		minimum height=0.72cm,
		inner sep=3pt,
		font=\sffamily\scriptsize
	},
	edge/.style={
		-{Stealth[length=2.1mm,width=1.5mm]},
		semithick,
		draw=dkgray!90
	},
	relation/.style={
		-{Stealth[length=2.1mm,width=1.5mm]},
		dashed,
		semithick,
		draw=cardinalred!65
	},
	label/.style={
		font=\sffamily\scriptsize,
		text=dkgray,
		fill=white,
		inner sep=1.5pt
	}
]

\node[root] (root) at (0.00,0.00) {Redaccion academica\\en entornos virtuales};

\node[branch] (modalidad) at (-5.82,1.42) {Educacion abierta\\y a distancia};
\node[branch] (autogestion) at (-1.77,3.81) {Autogestion del\\aprendizaje};
\node[branch] (comunicacion) at (1.77,3.82) {Comunicacion en\\entornos virtuales};
\node[branch] (criterios) at (5.84,1.42) {Criterios de\\calidad textual};
\node[branch] (responsabilidad) at (0.00,-3.35) {Responsabilidad\\academica};

\node[subbranch] (modalidad1) at (-8.85,3.15) {Flexibilidad\\espacial y temporal};
\node[subbranch] (modalidad2) at (-9.16,1.43) {Comunicacion\\mediada por texto};
\node[subbranch] (modalidad3) at (-8.85,-0.25) {Menor supervision\\inmediata};

\node[subbranch] (autogestion1) at (-2.35,5.25) {Organizacion\\del tiempo};
\node[subbranch] (autogestion2) at (-4.98,2.94) {Aprendizaje\\autonomo};

\node[subbranch] (comunicacion1) at (2.35,5.25) {Mensajes y\\participaciones};
\node[subbranch] (comunicacion2) at (4.97,2.91) {Productos\\escritos};

\node[subbranch] (criterios1) at (8.85,3.15) {Claridad};
\node[subbranch] (criterios2) at (9.15,1.45) {Coherencia};
\node[subbranch] (criterios3) at (8.85,-0.25) {Cohesion y\\formalidad};

\node[subbranch] (responsabilidad1) at (-3.75,-5.15) {Constancia y\\cumplimiento};
\node[subbranch] (responsabilidad2) at (0.00,-5.45) {Lectura, revision\\y correccion};
\node[subbranch] (responsabilidad3) at (3.75,-5.15) {Texto pertinente\\y riguroso};

\draw[edge] (root.west) -- node[label,above]{se configura en} (modalidad.east);
\draw[edge] (root.north west) -- node[label,left]{requiere} (autogestion.south);
\draw[edge] (root.north east) -- node[label,right]{se expresa en} (comunicacion.south);
\draw[edge] (root.east) -- node[label,above]{se valora por} (criterios.west);
\draw[edge] (root.south) -- node[label,right]{depende de} (responsabilidad.north);

\draw[edge] (modalidad.north west) -- (modalidad1.east);
\draw[edge] (modalidad.west) -- (modalidad2.east);
\draw[edge] (modalidad.south west) -- (modalidad3.east);

\draw[edge] (autogestion.north) -- node[label,left]{implica} (autogestion1.south);
\draw[edge] (autogestion.west) -- node[label,below]{favorece} (autogestion2.east);

\draw[edge] (comunicacion.north) -- (comunicacion1.south);
\draw[edge] (comunicacion.east) -- (comunicacion2.west);

\draw[edge] (criterios.north east) -- node[label,above]{incluyen} (criterios1.west);
\draw[edge] (criterios.east) -- (criterios2.west);
\draw[edge] (criterios.south east) -- (criterios3.west);

\draw[edge] (responsabilidad.south west) -- (responsabilidad1.north);
\draw[edge] (responsabilidad.south) -- (responsabilidad2.north);
\draw[edge] (responsabilidad.south east) -- (responsabilidad3.north);

\draw[relation] (modalidad2.north east) to[out=35,in=180]
	node[label,above,pos=0.45]{favorece}
	(comunicacion1.west);

\draw[relation] (autogestion2.south) to[out=-80,in=150]
	node[label,left,pos=0.45]{fortalece}
	(responsabilidad2.north west);

\draw[relation] (comunicacion2.south east) to[out=-25,in=165]
	node[label,below,pos=0.55]{se ajusta a}
	(criterios2.west);

\end{tikzpicture}
}%
\caption{Mapa conceptual sobre la relacion entre modalidad virtual, autogestion del aprendizaje y calidad de la redaccion academica.}
\label{fig:mapa-conceptual-redaccion-portafolio}
\end{figure}
\end{landscape}
\endgroup
\clearpage

\subsubsection{Conclusión}

El análisis permite afirmar que la redacción académica en entornos virtuales no debe entenderse como una técnica aislada de corrección, sino como el resultado visible de un sistema de trabajo. La modalidad educativa introduce flexibilidad, pero esa flexibilidad solo produce aprendizaje real cuando se convierte en organización, seguimiento, lectura atenta, revisión y responsabilidad sobre lo que se entrega. En ese sentido, escribir bien no es solo presentar un texto limpio; es demostrar control del proceso, capacidad de síntesis y criterio para comunicar conocimiento con precisión.

Desde mi perspectiva, la relación más importante que deja ver esta actividad es la siguiente: la autonomía no vale por sí sola si no produce textos claros, pertinentes y rigurosos. En una carrera como Derecho, donde la argumentación escrita organiza decisiones, interpretaciones y posturas, claridad, coherencia, cohesión y formalidad funcionan como indicadores concretos de disciplina intelectual. Por eso, la redacción académica en ambientes virtuales debe asumirse como una práctica de estructura, responsabilidad y juicio, y no como un requisito accesorio de presentación.


\clearpage
\subsection{Actividad 7. Analisis integrador del portafolio}

Como actividad de cierre, el portafolio no solo reune evidencias, sino que interpreta el proceso de aprendizaje que estas muestran en conjunto. El criterio de integracion consiste en demostrar que las actividades 1 a 6 no fueron ejercicios aislados, sino momentos sucesivos de una misma competencia: escribir con claridad, pertinencia y control del sentido en entornos virtuales.

La secuencia seleccionada responde a una logica pedagogica. Primero aparece el analisis de un conflicto comunicativo; despues, la identificacion de criterios de calidad textual; luego, la transferencia de esos criterios al ambito laboral; posteriormente, la evaluacion de un texto propio; mas adelante, la transformacion de un registro informal a uno academico; y, al final, la sintesis conceptual del sistema que sostiene la escritura en entornos virtuales. Esta organizacion permite observar continuidad y progresion, no solo acumulacion de productos.

\subsubsection{Integracion de actividades}

\begingroup
\small
\setlength{\tabcolsep}{5pt}
\renewcommand{\arraystretch}{1.28}
\begin{longtable}{@{}p{0.10\linewidth}p{0.23\linewidth}p{0.31\linewidth}p{0.28\linewidth}@{}}
\caption{Integracion de actividades 1 a 6}\label{tab:portafolio-evidencias}\\
\toprule
\textbf{Actividad} & \textbf{Producto} & \textbf{Problema trabajado} & \textbf{Aprendizaje recuperado} \\
\midrule
\endfirsthead
\toprule
\textbf{Actividad} & \textbf{Producto} & \textbf{Problema trabajado} & \textbf{Aprendizaje recuperado} \\
\midrule
\endhead
1 & Estudio de caso sobre proceso comunicativo & Analizar un conflicto entre companeros en una plataforma academica y distinguir emisor, receptor, mensaje, canal y contexto. & El desacuerdo solo es formativo cuando se expresa con pertinencia, respeto y finalidad constructiva. \\
\addlinespace
2 & Cuadro comparativo de criterios de redaccion academica & Contrastar un texto coloquial con otro academico para reconocer diferencias de claridad, coherencia, cohesion, formalidad y objetividad. & La eficacia academica depende de organizar la informacion y no solo de emitir una opinion. \\
\addlinespace
3 & Juego de roles con correos laborales & Reformular una solicitud urgente para ajustarla al contexto profesional mediante asunto claro, tono adecuado y estructura formal. & La redaccion profesional mejora la posibilidad de respuesta util cuando combina precision, cortesia y contexto. \\
\addlinespace
4 & Lista de verificacion para evaluar textos digitales & Revisar un texto propio para detectar fortalezas y limites de legibilidad en pantalla, estructura y tono. & Un texto correcto tambien debe ser visualmente legible y facil de recorrer en formato digital. \\
\addlinespace
5 & Ejercicio de transformacion de texto informal a academico & Convertir un parrafo coloquial sobre el estudio en linea en una version academica con conectores, precision y control discursivo. & Transformar el registro implica reorganizar el pensamiento y no solo sustituir palabras informales. \\
\addlinespace
6 & Mapa conceptual sobre autogestion del aprendizaje y redaccion academica & Relacionar modalidad virtual, autonomia, organizacion del tiempo, responsabilidad academica y calidad textual. & La escritura academica es resultado visible de un sistema de autogestion, disciplina y comunicacion. \\
\bottomrule
\end{longtable}
\endgroup

La integracion permite reconocer una progresion definida. Las actividades 1 y 2 trabajan el diagnostico: identifican el conflicto comunicativo y distinguen los criterios que separan una opinion espontanea de un texto academico funcional. Las actividades 3 y 4 trasladan esos criterios a contextos concretos de escritura profesional y evaluacion digital. Finalmente, las actividades 5 y 6 consolidan la apropiacion del aprendizaje, primero mediante la reescritura de un texto y luego mediante su organizacion conceptual.

Lo relevante es que las seis evidencias no se comportan como piezas aisladas. En conjunto, muestran que la comunicacion academica en entornos virtuales exige una misma logica de fondo: planificar lo que se dice, considerar al receptor, justificar el mensaje y revisar su forma final. La evidencia ya no es solo que se realizaron seis tareas, sino que se desarrollo una linea continua de aprendizaje sobre el uso responsable y estrategico de la escritura.

Conservar las seis actividades permite mostrar el proceso con mayor fidelidad que una seleccion limitada de productos. El portafolio cumple asi dos funciones complementarias: resguarda las evidencias elaboradas y ofrece una lectura comparativa que explica el aporte particular de cada una al desarrollo de la competencia comunicativa.

\section{Conclusion}

El portafolio de evidencias permite concluir que las actividades 1 a 6 no fueron ejercicios independientes, sino etapas de una misma formacion en comunicacion escrita. Cada producto atendio un problema particular, pero todos convergieron en una idea comun: escribir bien en entornos virtuales significa ordenar el pensamiento, adecuar el lenguaje al contexto y asumir responsabilidad sobre el efecto del mensaje.

Considero que mi principal avance fue comprender que la redaccion academica no se aprende por repeticion mecanica de reglas, sino mediante contraste, correccion y transferencia entre situaciones diferentes. La dificultad mas constante fue pasar de una expresion inmediata a una escritura capaz de conservar claridad, formalidad y sentido sin perder de vista al receptor. Revisar cada evidencia permitio reconocer que planificar, citar fuentes, reorganizar ideas y corregir la presentacion son partes de un mismo proceso.

Desde mi perspectiva, integrar las actividades completas ofrece una imagen fiel del proceso formativo porque permite observar tanto los avances como los aspectos que todavia requieren practica. En una licenciatura como Derecho, donde interpretar, argumentar, documentar y comunicar por escrito resulta indispensable, estas habilidades tienen utilidad academica y profesional directa. El valor final del portafolio no consiste solamente en reunir trabajos, sino en demostrar una forma mas consciente, responsable y estrategica de construir mensajes escritos.

% BIBLIOGRAFIA
% Agregar nuevas fuentes en redaccion-en-contextos-virtuales.bib.
% Citar en el texto con \citep{clave} o \citet{clave}.
\clearpage
\bibliography{redaccion-en-contextos-virtuales}

% FIN DEL DOCUMENTO
\end{document}


